% generated by roam.py handbook export -- не редактировать, править заметки
% 20260922165613-усиленное_неравенство_для_функции_с_липшицевым_градиентом.org (body 747b6a33a7b7)

\begin{theorem}[{Усиленное неравенство для функций с липшицевым градиентом}]\label[theorem]{note:379E4CAD-A44B-4631-8DAF-F3DC708EAEDE}
  Пусть $E \subset \R^n$~--- \hyperref[note:316B5BFA-8B9F-4267-9086-D6593AAB5F9B]{выпуклое множество}, $f: E \rightarrow \R$, $f \in C^{1,1}_L$~--- функция с \hyperref[note:7DD1C7AB-1A49-4675-9F21-5A16C717F670]{липшицевым градиентом} \\
  Тогда
  \[
  \forall x,y \in E\ \left| f(y) -f(x) - \left\langle \nabla f(x); y-x  \right\rangle \right|\le \frac{L}{2} \| y - x\|^2_2
  \]
\end{theorem}

\begin{remark}

  Если мы раскроем модуль и посмотрим на верхнюю границу, то получим:
  \[
  \forall x,y \in E\ f(y) - f(x) - \left\langle \nabla f(x); y-x \right\rangle \le \frac{L}{2} \| y -x \|^2
   \]
   Что равносильно $f(y) \le f(x) + \left\langle \nabla f(x); y-x \right\rangle + \frac{L}{2} \| y-x \|^2$,
   а это означает, что график функции лежит не выше параболы~--- двойственно \hyperref[note:22E546FC-186F-4A67-85AA-2C85DCEF9021]{сильной выпуклости}, где он лежит не ниже параболы
\end{remark}

\begin{proof}
  Заметим, что для функции $\phi \in C^1[0;1]$ верно (\hyperref[note:AD362344-ECDC-4627-B344-8356744E6E53]{формула Ньютона-Лейбница}), что $\phi(1) - \phi(0) = \int\limits_0^1 \phi^{\prime}(x)dx$ \\
  Давайте рассмотрим $\phi(t) = f(x + t(y-x))$, тогда $\phi(1) = f(y),\ \phi(0)=f(x)$ \\
  Тогда давайте представим $f(y)-f(x)$, как интеграл  от $\phi^{\prime}$.
  Мы знаем, что $\phi^{\prime}(t) = f^{\prime}(x + t(y-x)) (y-x) = \left\langle \nabla f(x + t(y-x)); y-x \right\rangle$

  Давайте это применим
  \begin{multline*}
  |f(y) -f(x) - \left\langle \nabla f(x); y-x  \right\rangle | =  \\
  \left| \int\limits_0^1 \left\langle \nabla f(x+t(y-x)); y-x \right\rangle dt - \left\langle \nabla f(x); y-x \right\rangle \right|
  \end{multline*}

  Так как $\left\langle \nabla f(x); y-x \right\rangle$~--- константа, ее можно представить так $\int\limits_0^1 \left\langle \nabla f(x); y-x \right\rangle dt$

  Давайте теперь сгруппируем интегралы, а подынтегральную сумму скалярных произведений
  схлопнем по линейности

  \begin{multline*}
  \left| \int\limits_0^1 \left\langle \nabla f(x+t(y-x)); y-x \right\rangle dt - \left\langle \nabla f(x); y-x \right\rangle \right| = \\
  \left| \int\limits_0^1 \left\langle \nabla f(x + t(y-x)) - \nabla f(x); y-x \right\rangle dt\right|
  \end{multline*}

  Еще мы знаем, что модуль интеграла не больше интеграла модуля

  \begin{multline*}
  \left| \int\limits_0^1 \left\langle \nabla f(x + t(y-x)) - \nabla f(x); y-x \right\rangle dt\right| \le
   \int\limits_0^1 \left| \left\langle \nabla f(x + t(y-x)) - \nabla f(x); y-x \right\rangle \right| dt
  \end{multline*}

  А также мы можем применить неравенство \hyperref[note:D0601ACD-8C91-4A93-B0E4-AC0D529C6A5C]{КБШ} для скалярного произведения
  \[
  \left| \left\langle \nabla f(x + t(y-x)) - \nabla f(x); y-x \right\rangle \right| \le \| \nabla f(x + t(y-x)) - \nabla f(x) \| \cdot \| y-x \|
  \]

  Если взять \hyperref[note:7DD1C7AB-1A49-4675-9F21-5A16C717F670]{условие липшицевости градиента} для точек $x + t(y-x)$ и $x$, то получим
  \[
  \| \nabla f(x +  t(y-x)) - \nabla f(x) \| \le L \| x + t(y-x) - x \| = t L \| y-x\|
  \]

  Итого
  \[
  \left| f(y) -f(x) - \left\langle \nabla f(x); y-x \right\rangle \right| \le \int\limits_{0}^1 t L \| y -x \|  \| y -x \| dt =  \frac{L}{2} \| y -x \|^2
  \]
\end{proof}
